\section{模型的评价与推广}

% 参考优秀论文 A196 的末章组织：评价必须对应前文已经得到的模型、结果和验证证据，
% 不是统一套用“优点很多、推广性强”等空泛表述。

\subsection{模型的优点}
% 数量由实际模型决定；每一点都说明“哪一项结构/机制/数值证据”带来了什么具体优势。
\begin{enumerate}
    \item ……。
    \item ……。
\end{enumerate}

\subsection{模型的不足与改进方向}
% 不足要说明会影响哪个结果、参数范围或适用边界；改进方向必须与该不足一一对应。
\begin{enumerate}
    \item ……。
    \item ……。
\end{enumerate}

\subsection{模型的推广}
% 选写。只有存在明确可迁移对象时保留，说明哪些模型结构可直接继承、哪些参数/边界需要重新标定。

% 中文国赛默认不单设“结论”一级章；各问题的最终回答应已经在对应“求解结果/结果分析”中闭环。
